\sect{Обзор существующих работ}{related}
%
Существующие работы по отслеживанию партитуры удобно сгруппировать
в четыре тематических направления, отражающих историческую
последовательность развития области и одновременно различающихся
по типу применяемого математического аппарата. Ниже последовательно
рассматриваются: классические алгоритмы динамического
программирования и их каузальные модификации; вероятностные
модели на основе скрытых марковских процессов и их полу-марковских
расширений; современные нейросетевые подходы; и относительно
недавно сформировавшееся семейство методов, в которых отслеживание
партитуры формулируется как задача обучения с подкреплением.

\textbf{Классические алгоритмы динамического программирования.}
Современные методы выравнивания временных рядов берут начало с
работы Сакоэ и Чибы~\cite{SakoeChiba1978}, в которой для задачи
распознавания речи введён алгоритм Dynamic Time Warping и его
теперь стандартная полоса ограничений. Перенос идей в задачу
отслеживания партитуры независимо предложен Данненбергом и
Веркоу~\cite{Dannenberg1984,Vercoe1984}: оба автора рассматривают
исполнение и партитуру как два потока однотипных событий и
сопоставляют их по принципу динамического программирования с
допуском ошибок. Существенным каузальным расширением
является On-line Time Warping
Диксона~\cite{Dixon2005}, в котором матрица накопленных
стоимостей заполняется столбец за столбцом по мере поступления
новых наблюдений, а специальный счётчик \textit{RunCount}
ограничивает число последовательных шагов в одном направлении и
тем самым предотвращает «застревание» алгоритма. Подробное
изложение DTW и музыкальных применений приведено в
монографии Мюллера~\cite{Mueller2007} и более поздней
учебной книге~\cite{MullerFMP2015}, а в диссертации
Арцта~\cite{Arzt2016} систематизирован опыт построения
надёжных DP-следователей под концертные условия.

\textbf{Вероятностные модели на основе скрытых марковских
процессов.} Учебная статья Рабинера~\cite{Rabiner1989} остаётся
канонической ссылкой для алгоритмов Forward, Backward, Viterbi
и переоценки параметров методом Баума-Уэлча; именно эта нотация
используется во всём дальнейшем изложении. Применение HMM к
задаче отслеживания партитуры впервые систематически проведено
в работе Орио и Дешелля~\cite{OrioDechelle2001}: предложена
двухуровневая модель, в которой нижний уровень моделирует
ноты как фазы attack--sustain--silence, а верхний~--- допустимые
ошибки исполнителя через параллельные «теневые» состояния. Конт
с соавторами~\cite{Cont2006} развил эту линию до иерархической
HMM с разреженными неотрицательными ограничениями для
полифонического выравнивания. Ключевым теоретическим достижением
оказалась работа Конта~\cite{Cont2010}, в которой предложена
\textit{гибридная} архитектура из двух связанных вероятностных
агентов~--- аудио-агента и темпо-агента~--- работающих в одной
байесовской инференц-схеме; при этом темпоральные элементы
партитуры моделируются как полу-марковский процесс с явным
распределением длительностей $d_j(u)$, а атемпоральные
(орнаменты, трели)~--- как обычный марковский процесс с
геометрической длительностью. Эта архитектура положена в основу
системы Antescofo, активно используемой в концертной практике
IRCAM с 2007 года. Параллельно Рафаэль~\cite{Raphael2010}
развивал родственный, но самостоятельный подход~---
переключающиеся пространства состояний с явным гауссовским
тайминг-моделем для пары $(t_n,s_n)$, объединённой с
HMM-наблюдательной частью; в системе Music Plus One этот подход
позволяет адаптироваться к конкретному солисту от репетиции к
репетиции через off-line EM. Накамура с соавторами~\cite{Nakamura2015}
рассмотрели задачу в наиболее общей постановке~--- с
произвольными повторами, пропусками и фальшивыми нотами~--- и
предложили факторизацию переходного графа, понижающую сложность
шага Forward с $O(N^2)$ до $O(N)$.

\textbf{Современные нейросетевые подходы.} С 2018 года активно
развивается семейство методов, обучаемых сквозным образом на
больших корпусах выровненных аудиозаписей. Хоторн с
соавторами~\cite{Hawthorne2018} в работе Onsets and Frames
ввели двойную целевую функцию~--- покадровое присутствие ноты
плюс отдельную голову для onset-событий; их же работа
2019~года~\cite{Hawthorne2019Maestro} ввела датасет MAESTRO
объёмом $172.3$~ч с точностью выравнивания около $3$\,мс,
который стал стандартным бенчмарком для пианистической
транскрипции. Арцт и Латтнер~\cite{ArztLattner2018} предложили
обучаемые транспозиционно-инвариантные признаки на основе
gated-autoencoder’а, превратив задачу audio-to-score в задачу
audio-to-audio выравнивания после синтеза партитуры через
семплер. Хенкель и Видмер~\cite{Henkel2021} разработали
систему Conditional YOLO с условием FiLM-параметров от аудио,
ведущую следование непосредственно по \textit{изображению}
партитуры на трёх разрешениях (нота, такт, система); такая
постановка снимает зависимость от дорогой нотно-уровневой
разметки. Аграваль и Диксон~\cite{AgrawalDixon2019} рассмотрели
гибридную схему, в которой обученный сиамской CNN признак
заменяет ручную chroma в матрице сходства, а само выравнивание
по-прежнему ведётся методом DTW.

\textbf{Обучение с подкреплением для отслеживания партитуры и
аккомпанемента.} Принципиально иная постановка предложена
Дорфером и коллегами~\cite{Dorfer2018RL}: задача формулируется
как мультимодальный марковский процесс принятия решений, в
котором состояние объединяет окно изображения партитуры с
участком спектрограммы, а действие изменяет скорость продвижения
курсора по партитуре с шагом $\Delta v_{\text{pxl}}$. Обучение
ведётся методом advantage actor--critic с шестнадцатью
параллельными акторами и плотной линейной функцией
вознаграждения~$r=1-|d_x|/b$ внутри окна допустимых отклонений.
Петер~\cite{Peter2023} предложил вариант для символьного
выравнивания с офлайновым обучением через offline RL: трансформер
аппроксимирует Q-функцию, но с дисконтом $\gamma=0$ (полностью
близорукий агент), благодаря чему задача сводится к бинарной
классификации правильности выбранного onset-кандидата. Его
система достигает медианной асинхронности $15.7$\,мс
(против $60.6$\,мс у OLTW), что является лучшей известной
авторам цифрой для онлайн-выравнивания символьного входа.
Цзян с соавторами~\cite{Jiang2020RLDuet} в системе RL-Duet
применили актор-критик с обобщённой оценкой
преимуществ~\cite{Schulman2016GAE} к задаче генерации
гармонического сопровождения для монофонической мелодии; функция
вознаграждения собирается из ансамбля
четырёх специализированных нейросетевых критиков. Ву с
соавторами~\cite{Wu2024ReaLchords} в работе ReaLchords развили
эту линию через KL-контроль к офлайновому учителю, видящему всю
мелодию целиком; такой подход формально решает проблему
\textit{exposure bias} при онлайн-генерации. Жак с
соавторами~\cite{Jaques2017SeqTutor} ранее предложили
KL-контроль как общий приём дообучения предобученной MLE-модели
последовательности с эвристическим внешним вознаграждением.
Обзорная работа Кима с соавторами~\cite{Kim2026LiveAgents}
систематизирует $184$ систему живых музыкальных агентов в
четырёх измерениях (Usage Context, Interaction, Technology,
Ecosystem) с $31$ детальной осью.

\textbf{Перенос методов из астроинформатики.} Отдельным
направлением, ранее не представленным в литературе по MIR,
являются методы slotted symbolic Markov modeling
(SSMM)~\cite{JohnstonPeter2017} и distribution
fields~\cite{SevillaLaraDF2012}, изначально развитые для
классификации переменных звёзд по нерегулярным фотометрическим
рядам~\cite{Johnston2020}. Перенос этих методов в задачу
отслеживания партитуры обсуждается в~\Secref{ssmm-df}.

\textbf{Положение настоящей работы.} В отличие от
работ~\cite{Dorfer2018RL,Peter2023}, в которых RL-агент
\textit{заменяет} классический трекер целиком, и работ
\cite{Jiang2020RLDuet,Wu2024ReaLchords}, в которых RL применяется
к генерации музыки, а не к её отслеживанию, в~\Secref{proposal}
обсуждается теоретическая возможность построения гибридной
схемы: классический вероятностный трекер сохраняется как ядро,
а RL-агент действует как \textit{тонкий слой опережающего
предсказания темпа} с целью компенсации задержек воспроизводящей
части системы. Такая постановка не противоречит результатам
указанных работ и формально дополняет их.
